\documentclass[10pt,aspectratio=169]{beamer}
\usetheme{default}
\usefonttheme{professionalfonts}
\setbeamertemplate{navigation symbols}{}
\usepackage[T1]{fontenc}
\usepackage{lmodern}
\usepackage{amsmath}
\usepackage{booktabs}
\usepackage{array}
\usepackage{hyperref}

\definecolor{ink}{HTML}{1A2744}
\definecolor{accent}{HTML}{1F4E79}
\definecolor{mute}{HTML}{5A6570}
\definecolor{card}{HTML}{F4F6F8}

\setbeamercolor{normal text}{fg=ink,bg=white}
\setbeamercolor{frametitle}{fg=ink,bg=white}
\setbeamercolor{title}{fg=ink,bg=white}
\setbeamercolor{block title}{fg=white,bg=accent}
\setbeamercolor{block body}{bg=card,fg=ink}
\setbeamercolor{itemize item}{fg=accent}
\setbeamerfont{title}{series=\bfseries,size=\LARGE}
\setbeamerfont{frametitle}{series=\bfseries,size=\large}

\setbeamertemplate{footline}{%
  \leavevmode%
  \hbox{%
    \begin{beamercolorbox}[wd=.82\paperwidth,ht=2.4ex,dp=1ex,leftskip=1.2em]{author in head/foot}%
      {\color{mute}\scriptsize Numerical Analysis --- Project Proposal}%
    \end{beamercolorbox}%
    \begin{beamercolorbox}[wd=.18\paperwidth,ht=2.4ex,dp=1ex,center]{date in head/foot}%
      {\color{mute}\scriptsize\insertframenumber/\inserttotalframenumber}%
    \end{beamercolorbox}%
  }%
}

\title{Comparing Three Numerical Methods\\for Planar Two-Link Arm Motion}
\date{}

\begin{document}

%----------------------------------------------------------------------
\begin{frame}[t]
  \vspace{0.15cm}
  \noindent{\LARGE\bfseries Comparing Three Numerical Methods\\[0.15em]
  for Planar Two-Link Arm Motion}

  \vspace{0.45cm}
  {\small
  \textbf{Course.} CSE402: Numerical Analysis, Simulation and Modeling Sessional\\[0.25em]
  \textbf{Section.} C-2\\[0.4em]
  \textbf{Group members.}
  \begin{itemize}
    \item 2105161 --- Sadia Naushin
    \item 2105163 --- Dibbo Chowdhury
    \item 2105173 --- Sunanda Chowdhury
    \item 2105174 --- Nujhat Sadia
    \item 2105176 --- Navin Nawar
    \item 2105179 --- Arpa Das
  \end{itemize}
  }
\end{frame}

%----------------------------------------------------------------------
\begin{frame}{Problem Statement \& Scope}
  \textbf{Problem.} Compute $q(t),\tau(t)$ on a fixed horizon $[0,T]$ taking a planar two-link arm from pose~$A$ to pose~$B$.
  With $q=(\theta_1,\theta_2)^\top$ and $\tau=(\tau_1,\tau_2)^\top$,
  {\small
  \begin{align*}
    M(q)\,\ddot q + C(q,\dot q)\,\dot q + g(q) &= \tau, && t\in[0,T],\\
    q(0)=q_A,\quad q(T)&=q_B,\\
    q_{\min} &\le q(t) \le q_{\max}.
  \end{align*}
  }

  \vspace{0.15cm}
  \textbf{Approach (course directions).}
  \begin{itemize}
    \item \textbf{Cross-domain application:} course numerical tools (time integration, collocation, solving nonlinear systems) applied to robot-arm motion.
    \item \textbf{Exploratory / comparative study:} the same model and task, three established methods, compared on accuracy and cost. Not a new algorithm.
  \end{itemize}

  \vspace{0.15cm}
  \begin{block}{Scope}
    {\small Simulation only (Python). Two-link planar arm. We do not reproduce the base paper's high-DoF pipeline or use hardware.}
  \end{block}
\end{frame}

%----------------------------------------------------------------------
\begin{frame}{Base Paper Overview}
  \textbf{Cardona-Ortiz, D.\ and Arechavaleta, G.} (2025).\\
  Trajectory optimization for highly articulated robots based on sparsity-free local direct collocation.\\
  \emph{International Journal of Applied Mathematics and Computer Science}, 35(4), 577--589.\\
  \url{https://doi.org/10.61822/amcs-2025-0041}

  \vspace{0.4cm}
  \textbf{Domain.} Numerical optimal control for robot motion (many-joint arms and similar mechanisms).

  \vspace{0.3cm}
  \textbf{Summary.} The paper plans robot motions by direct collocation: trapezoidal and Hermite--Simpson rules turn continuous physics into a finite set of equations, which are then solved numerically. Their extra software (sparse Jacobian assembly, Pinocchio, high-DoF platforms) is not our implementation target.

  \vspace{0.3cm}
  {\small Method recipes we follow in code are those in Kelly, \emph{An Introduction to Trajectory Optimization}, SIAM Review, 2017.}
\end{frame}

%----------------------------------------------------------------------
\begin{frame}{Expected Methodology}
  \textbf{Model.} Planar two-link arm; states are joint angles and velocities; inputs are joint torques.

  \vspace{0.3cm}
  \textbf{Methods (identical start, goal, and time budget).}
  \begin{enumerate}
    \item Shooting with fourth-order Runge--Kutta integration
    \item Trapezoidal direct collocation
    \item Hermite--Simpson direct collocation
  \end{enumerate}

  \vspace{0.25cm}
  \textbf{Evaluation.} Tip error at the goal, residual of the dynamics, and runtime, versus number of time nodes.\\[0.35em]
  \textbf{Deliverables.} Implementation; comparison table and plots; synchronized matplotlib visualization of the three motions.
\end{frame}

\end{document}
